% ============================================================
%  BEAMER — MATEMÁTICA II · GUÍA DE EJERCICIOS Nº 1
%  Funciones reales de varias variables
%  Incrementos parciales y totales
%  Autor: Ing. Gabriel Astudillo
%  Compilar:  latexmk -xelatex beamer_mat2_guia1.tex
% ============================================================
\documentclass[aspectratio=169]{beamer}
\usepackage{mathwizards-beamer}
\mwbeamer{Matemática II}{Guía de Ejercicios Nº 1}%
  {Funciones de varias variables · Incrementos parciales y totales}

\begin{document}

% ════════════════════════════════════════════════════════════
\begin{frame}[plain]
  \titlepage
\end{frame}

% ════════════════════════════════════════════════════════════
\begin{frame}{Fuente y atribución}
  \begin{mwextra}[Material original]
    Díaz L., Isabel. \emph{Guía de Ejercicios Nº 1 — Matemática II}.
    Universidad Central de Venezuela (UCV), enero de 2017.
  \end{mwextra}
  \begin{block}{Esta presentación}
    Selección de ejercicios, desarrollo teórico-práctico y adaptación
    al estilo Mathwizards: Ing. Gabriel Astudillo
    (Mathwizards Consultoría Educativa STEM).
  \end{block}
  \begin{alertblock}{Uso del material}
    Material didáctico elaborado con autorización de la fuente original.
    Se cita la autoría y la institución de procedencia.
  \end{alertblock}
\end{frame}

% ════════════════════════════════════════════════════════════
\begin{frame}{Objetivos y mapa de la guía}
  \begin{block}{Al terminar la clase podrás}
    \begin{itemize}
      \item identificar, evaluar e interpretar funciones reales de varias variables;
      \item construir funciones económicas de demanda, costo, ingreso y beneficio;
      \item reconocer funciones homogéneas y funciones de producción Cobb-Douglas;
      \item calcular e interpretar incrementos parciales y totales.
    \end{itemize}
  \end{block}
  \vspace{2pt}
  \begin{columns}[T]
    \column{0.48\textwidth}
    \begin{mwextra}[Sección 1 de la guía]
      Función real de varias variables independientes.\\
      Ejercicios \textbf{5, 10, 11, 17, 20, 21 y 22}.
    \end{mwextra}
    \column{0.48\textwidth}
    \begin{mwextra}[Sección 2 de la guía]
      Incrementos parciales y totales.\\
      Ejercicios \textbf{26, 27, 28, 31, 33 y 34}.
    \end{mwextra}
  \end{columns}
\end{frame}

% ════════════════════════════════════════════════════════════
\section{Marco Teórico}
% ════════════════════════════════════════════════════════════

\begin{frame}{Función real de varias variables}
  \begin{block}{Definición}
    Una \textbf{función real de \(n\) variables reales} es una regla que asigna
    a cada punto \(\mathbf{x}=(x_1,\dots,x_n)\) de un conjunto
    \(D\subseteq\mathbb{R}^{n}\) un único valor real \(f(\mathbf{x})\):
    \[
      f:D\subseteq\mathbb{R}^{n}\longrightarrow\mathbb{R},
      \qquad \mathbf{x}\longmapsto f(\mathbf{x}).
    \]
    \(D\) es el \textbf{dominio} y el conjunto de valores que toma \(f\) es el
    \textbf{rango}.
  \end{block}
  \begin{columns}[T]
    \column{0.48\textwidth}
    \begin{mwextra}[Dos variables]
      \(z=f(x,y)\): el dominio es una región del plano y la gráfica es una
      \textbf{superficie} en \(\mathbb{R}^{3}\).
    \end{mwextra}
    \column{0.48\textwidth}
    \begin{mwextra}[Tres variables]
      \(w=f(x,y,z)\): el dominio es una región de \(\mathbb{R}^{3}\); la gráfica
      vive en \(\mathbb{R}^{4}\) y no se representa directamente.
    \end{mwextra}
  \end{columns}
\end{frame}

\begin{frame}{Evaluar una función de varias variables}
  \begin{mwpasos}[Procedimiento]
    \begin{enumerate}
      \item Sustituye cada variable por su valor, respetando el orden del argumento.
      \item Opera con las reglas de siempre: potencias, raíces, logaritmos.
      \item Interpreta el resultado con sus unidades.
    \end{enumerate}
  \end{mwpasos}
  \vspace{2pt}
  \begin{columns}[T]
    \column{0.5\textwidth}
    \begin{exampleblock}{Ejemplo}
      Si \(f(x,y)=x^{2}y-3y+5\),
      \[
        f(2,3)=2^{2}(3)-3(3)+5=8.
      \]
    \end{exampleblock}
    \column{0.46\textwidth}
    \begin{alertblock}{Cuidado}
      \(f(x+\Delta x)\neq f(x)+\Delta x\). Primero se sustituye \emph{todo} el
      argumento y después se opera.
    \end{alertblock}
  \end{columns}
\end{frame}

\begin{frame}{Dominio y rango: qué revisar}
  \begin{block}{Restricciones algebraicas}
    \begin{itemize}
      \item Denominadores: \(D(x)\neq 0\).
      \item Raíces de índice par: radicando \(\ge 0\).
      \item Logaritmos: argumento \(>0\).
    \end{itemize}
  \end{block}
  \begin{mwextra}[Restricciones del contexto económico]
    Precios, cantidades, capital y trabajo suelen exigir \(p,q,L,K\ge 0\).
    El dominio matemático puede ser mayor que el dominio con sentido económico.
  \end{mwextra}
  \begin{exampleblock}{Ejemplo del ejercicio 28}
    \(C(x,y,t)=(100+200x+10y)e^{0{,}01t}\) es un polinomio por una exponencial:
    su dominio matemático es \(\mathbb{R}^{3}\). En el contexto, \(x,y\ge 0\)
    y \(t\ge 0\).
  \end{exampleblock}
\end{frame}

\begin{frame}{Modelos económicos: variables y relaciones}
  \begin{columns}[T]
    \column{0.48\textwidth}
    \begin{block}{Tipos de variables}
      \begin{itemize}
        \item \textbf{Endógenas:} las determina el modelo
              (\(q,I,C,B\)).
        \item \textbf{Exógenas:} vienen dadas
              (\(p\) en competencia perfecta, \(S\), \(t\)).
      \end{itemize}
    \end{block}
    \column{0.48\textwidth}
    \begin{block}{Mercado}
      Demanda: \(q_d=D(p)\)\\[2pt]
      Oferta: \(q_s=S(p)\)\\[2pt]
      Equilibrio: \(D(p)=S(p)\)
    \end{block}
  \end{columns}
  \vspace{2pt}
  \begin{mwpasos}[Del enunciado a la función]
    \begin{enumerate}
      \item Identifica la variable que se quiere explicar.
      \item Identifica las variables de las que depende y sus unidades.
      \item Escribe las relaciones económicas y sustituye.
      \item Simplifica y declara el dominio con sentido económico.
    \end{enumerate}
  \end{mwpasos}
\end{frame}

\begin{frame}{Ingreso, costo y beneficio}
  \begin{block}{Identidades básicas}
    \[
      I=pq,\qquad C=C_F+C_V(q),\qquad B(q)=I(q)-C(q).
    \]
  \end{block}
  \begin{columns}[T]
    \column{0.48\textwidth}
    \begin{mwextra}[Monopolio]
      El precio depende de la cantidad: \(p=p(q)\). Por eso
      \[
        I(q)=p(q)\,q.
      \]
    \end{mwextra}
    \column{0.48\textwidth}
    \begin{mwextra}[Competencia perfecta]
      El precio lo fija el mercado: \(p=p^{*}\) constante. Entonces
      \[
        I(q)=p^{*}q.
      \]
    \end{mwextra}
  \end{columns}
  \vspace{2pt}
  \begin{alertblock}{Costo conjunto}
    Si se producen dos bienes: \(C=C(q_X,q_Y)\). Si hay dos plantas:
    \(C=C_X(q_X)+C_Y(q_Y)\), con \(q=q_X+q_Y\).
  \end{alertblock}
\end{frame}

\begin{frame}{Varios productos, mercados o plantas}
  \begin{block}{Producción múltiple (ejercicio 11A)}
    \[
      B(q_X,q_Y)=I(q_X,q_Y)-C(q_X,q_Y).
    \]
    En monopolio, \(I=p_X(q_X,q_Y)\,q_X+p_Y(q_X,q_Y)\,q_Y\). En competencia
    perfecta, \(p_X=p_X^{*}\) y \(p_Y=p_Y^{*}\) son constantes.
  \end{block}
  \begin{block}{Discriminación de precios (ejercicio 11B)}
    \[
      B(q_X,q_Y)=I_X(q_X)+I_Y(q_Y)-C(q_X+q_Y).
    \]
  \end{block}
  \begin{block}{Dos fábricas, un solo producto (ejercicio 11C)}
    \[
      B(q_X,q_Y)=I(q_X+q_Y)-C_X(q_X)-C_Y(q_Y).
    \]
  \end{block}
\end{frame}

\begin{frame}{Cobb-Douglas y funciones homogéneas}
  \begin{block}{Función de producción Cobb-Douglas}
    \[
      q(L,K)=A\,L^{\alpha}K^{1-\alpha},\qquad A>0,\quad 0<\alpha<1.
    \]
    \(L\): trabajo;\quad \(K\): capital;\quad \(q\): producción.
  \end{block}
  \begin{block}{Homogeneidad de grado \(n\)}
    \(f\) es homogénea de grado \(n\) si, para todo \(\lambda>0\),
    \[
      f(\lambda x,\lambda y)=\lambda^{n}f(x,y).
    \]
  \end{block}
  \begin{mwextra}[Lectura económica]
    En Cobb-Douglas, \(q(\lambda L,\lambda K)=\lambda\,q(L,K)\): si se duplican
    trabajo y capital, la producción se duplica. Son \textbf{rendimientos
    constantes a escala} (grado 1).
  \end{mwextra}
\end{frame}

\begin{frame}{Demanda lineal a partir de dos datos}
  \small\vspace{-1pt}
  \begin{block}{Modelo lineal}
    \[
      q(p)=a+bp,\qquad b=\frac{\Delta q}{\Delta p}<0.
    \]
    \(a\): demanda cuando \(p=0\);\quad \(b\): cambio de \(q\) por cada unidad
    de \(p\).
  \end{block}
  \begin{mwpasos}[Cómo construirla]
    \begin{enumerate}
      \item Con dos pares \((p_1,q_1)\) y \((p_2,q_2)\), calcula
            \(b=\dfrac{q_2-q_1}{p_2-p_1}\).
      \item Sustituye un par para despejar \(a=q_1-bp_1\).
      \item Escribe \(q(p)=a+bp\) y verifica con el otro par.
    \end{enumerate}
  \end{mwpasos}
  \begin{exampleblock}{Ejemplo genérico}
    Si \(q(10)=100\) y \(q(12)=90\), entonces \(q(p)=150-5p\).
  \end{exampleblock}
\end{frame}

\begin{frame}{Composición y sustitución de variables}
  \begin{block}{Idea}
    Si una magnitud depende de variables que a su vez dependen de otra, se
    sustituye para obtener una función de una sola variable.
  \end{block}
  \[
    B=B(P,A),\qquad P=P(T),\qquad A=A(T)
    \quad\Longrightarrow\quad
    B(T)=B\bigl(P(T),A(T)\bigr).
  \]
  \begin{mwpasos}[Orden de trabajo]
    \begin{enumerate}
      \item Escribe la función original con sus variables.
      \item Sustituye cada variable por su expresión en \(T\).
      \item Simplifica y declara el dominio (por ejemplo, \(T\ge 1\) si
            \(A=T-1\)).
    \end{enumerate}
  \end{mwpasos}
\end{frame}

\begin{frame}{Distancias y triángulos rectángulos}
  \begin{block}{Distancia entre dos puntos}
    \[
      d\bigl((x_1,y_1),(x_2,y_2)\bigr)
      =\sqrt{(x_2-x_1)^{2}+(y_2-y_1)^{2}}.
    \]
  \end{block}
  \begin{block}{Teorema de Pitágoras}
    En un triángulo rectángulo de catetos \(a\) y \(b\) e hipotenusa \(c\):
    \[
      c^{2}=a^{2}+b^{2}
      \quad\Longrightarrow\quad
      c=\sqrt{a^{2}+b^{2}}.
    \]
  \end{block}
  \begin{exampleblock}{Ejemplo}
    La distancia entre \((0,7)\) y \((4,4)\) es
    \(\sqrt{4^{2}+3^{2}}=\sqrt{25}=5\). Cada tramo diagonal de los ejercicios
    10 y 21 se calcula así.
  \end{exampleblock}
\end{frame}

\begin{frame}{Incrementos parciales}
  \begin{block}{Definición}
    El \textbf{incremento parcial} mide el cambio de \(f\) cuando varía
    \emph{una sola} variable y las demás permanecen fijas:
    \[
      \Delta_x f(x_0,y_0)=f(x_0+\Delta x,\,y_0)-f(x_0,y_0),
    \]
    \[
      \Delta_y f(x_0,y_0)=f(x_0,\,y_0+\Delta y)-f(x_0,y_0).
    \]
    En tres variables se define análogamente \(\Delta_z f\).
  \end{block}
  \begin{mwextra}[Interpretación]
    Es un cambio \emph{ceteris paribus}: «todo lo demás constante». El signo
    indica si la función aumenta o disminuye.
  \end{mwextra}
\end{frame}

\begin{frame}{Incremento total}
  \begin{block}{Definición}
    El \textbf{incremento total} permite que cambien \emph{varias} variables
    a la vez:
    \[
      \Delta f(x_0,y_0)=f(x_0+\Delta x,\,y_0+\Delta y)-f(x_0,y_0).
    \]
    En tres variables:
    \[
      \Delta f=f(x_0+\Delta x,\,y_0+\Delta y,\,z_0+\Delta z)
             -f(x_0,y_0,z_0).
    \]
  \end{block}
  \begin{alertblock}{No es simplemente la suma de los parciales}
    \[
      \Delta f=\Delta_x f(x_0,y_0)+\Delta_y f(x_0+\Delta x,\,y_0).
    \]
    En general aparecen \textbf{términos cruzados} (\(\Delta x\,\Delta y\)).
    Solo si \(f\) es lineal, \(\Delta f=\Delta_x f+\Delta_y f\).
  \end{alertblock}
\end{frame}

\begin{frame}{Incrementos en funciones cuadráticas}
  \begin{block}{Desarrollos útiles}
    \[
      (x+\Delta x)^{2}=x^{2}+2x\,\Delta x+(\Delta x)^{2},
    \]
    \[
      (x+\Delta x)(y+\Delta y)
      =xy+x\,\Delta y+y\,\Delta x+\Delta x\,\Delta y.
    \]
  \end{block}
  \begin{exampleblock}{Caso general \(f(x,y)=ax^{2}+by^{2}+cxy\)}
    \[
      \Delta_x f=(2ax+cy)\Delta x+a(\Delta x)^{2},
    \]
    \[
      \Delta_y f=(2by+cx)\Delta y+b(\Delta y)^{2},
    \]
    \[
      \Delta f=\Delta_x f+\Delta_y f+c\,\Delta x\,\Delta y.
    \]
  \end{exampleblock}
\end{frame}

\begin{frame}{Interpretar un incremento}
  \begin{columns}[T]
    \column{0.5\textwidth}
    \begin{block}{Qué reportar}
      \begin{itemize}
        \item \textbf{Signo:} aumenta (\(>0\)) o disminuye (\(<0\)).
        \item \textbf{Magnitud:} cuánto cambia, con unidades.
        \item \textbf{Contexto:} ¿conviene o no la decisión?
      \end{itemize}
    \end{block}
    \column{0.46\textwidth}
    \begin{mwextra}[Cambio porcentual]
      \[
        \%\,\Delta f\approx \frac{\Delta f}{f_0}\times 100.
      \]
      Útil cuando se compara un año con otro.
    \end{mwextra}
  \end{columns}
  \vspace{2pt}
  \begin{alertblock}{Errores frecuentes}
    Mezclar unidades (miles vs. unidades), olvidar el signo de \(\Delta\), o
    sumar incrementos parciales cuando el cambio es simultáneo.
  \end{alertblock}
\end{frame}

\begin{frame}{Estrategia para resolver los ejercicios}
  \small\vspace{-3pt}
  \begin{mwpasos}[Checklist]
    \begin{enumerate}
      \item \textbf{Variable dependiente:} ¿qué se quiere calcular?
      \item \textbf{Variables independientes:} ¿de qué depende y en qué unidades?
      \item \textbf{Datos:} precios, cantidades, costos fijos, puntos de demanda.
      \item \textbf{Modelo:} escribe las identidades y sustituye.
      \item \textbf{Cálculo:} evalúa o incrementa según lo pedido.
      \item \textbf{Interpretación:} traduce el número a una frase económica.
    \end{enumerate}
  \end{mwpasos}
  \begin{columns}[T]
    \column{0.48\textwidth}
    \begin{mwextra}[Si piden «demuestra»]
      Desarrolla un lado de la igualdad hasta obtener el otro; no sustituyas
      valores numéricos salvo que el enunciado lo pida.
    \end{mwextra}
    \column{0.48\textwidth}
    \begin{mwextra}[Si piden «interpreta»]
      Responde con unidades y una conclusión.
    \end{mwextra}
  \end{columns}
\end{frame}

\begin{frame}{Formulario de bolsillo}
  \begin{columns}[T]
    \column{0.48\textwidth}
    \begin{block}{Economía}
      \begin{itemize}
        \item \(I=pq\)
        \item \(B=I-C\)
        \item \(C=C_F+cq\)
        \item \(q_d=D(p)\), \(q_s=S(p)\)
        \item Equilibrio: \(D(p)=S(p)\)
      \end{itemize}
    \end{block}
    \column{0.48\textwidth}
    \begin{block}{Incrementos}
      \begin{itemize}
        \item \(\Delta_x f=f(x+\Delta x,y)-f(x,y)\)
        \item \(\Delta_y f=f(x,y+\Delta y)-f(x,y)\)
        \item \(\Delta f=f(x+\Delta x,y+\Delta y)-f(x,y)\)
        \item \(\%\Delta f\approx \dfrac{\Delta f}{f_0}\,100\)
      \end{itemize}
    \end{block}
  \end{columns}
  \vspace{2pt}
  \begin{mwextra}[Cobb-Douglas]
    \(q=AL^{\alpha}K^{1-\alpha}\) es homogénea de grado 1:
    \(q(\lambda L,\lambda K)=\lambda\,q(L,K)\).
  \end{mwextra}
\end{frame}

% ════════════════════════════════════════════════════════════
\section{Ejercicios}
% ════════════════════════════════════════════════════════════

\begin{frame}{Ejercicios seleccionados}
  \begin{block}{Guía de Ejercicios Nº 1 — MAT II}
    \begin{columns}[T]
      \column{0.48\textwidth}
      \textbf{Sección 1: función de varias variables}\\[2pt]
      5 · 10 · 11 · 17 · 20 · 21 · 22
      \column{0.48\textwidth}
      \textbf{Sección 2: incrementos parciales y totales}\\[2pt]
      26 · 27 · 28 · 31 · 33 · 34
    \end{columns}
  \end{block}
  \vspace{2pt}
  \begin{mwextra}[Cómo usar estas diapositivas]
    Cada enunciado queda anclado arriba; el espacio libre de abajo es para
    resolver a mano sobre el PDF.
  \end{mwextra}
\end{frame}

% ── Ejercicio 5 ─────────────────────────────────────────────
\begin{frame}{Ejercicio 5}
  \begin{mwejercicio}[Ejercicio 5 — Cobb-Douglas \quad \medio]
    La producción \(q\) de una fábrica se considera función de la inversión de
    capital \(K\) y del tamaño de la fuerza laboral \(L\). Las funciones
    \[
      q(L,K)=A\,L^{\alpha}K^{1-\alpha},\qquad A>0,\quad 0<\alpha<1,
    \]
    se conocen como \textbf{funciones de producción de Cobb-Douglas}.
    Supóngase que \(A=10\) y \(\alpha=\tfrac12\).

    Demuestra que la función de producción de esta fábrica satisface
    \[
      q(\alpha L,\alpha K)=\alpha\,q(L,K)
    \]
    y luego explica su significado.
  \end{mwejercicio}
\end{frame}

% ── Ejercicio 10 ────────────────────────────────────────────
\begin{frame}{Ejercicio 10}
  \mwtopanchoroff
  \begin{columns}[T]
    \column{0.54\textwidth}
    \begin{mwejercicio}[Ejercicio 10 — Conducto \quad \dificil]
      \small
      Se debe construir un conducto de agua desde \(P\) hasta \(S\). El
      conducto atravesará regiones con diferentes costos de construcción. El
      costo neto por kilómetro, en miles de dólares, es de \(300\) mil de
      \(P\) a \(Q\), \(200\) mil de \(Q\) a \(R\) y \(100\) mil de \(R\) a
      \(S\).
      \begin{itemize}
        \item[\textbf{(A)}] Determina analíticamente una expresión para el costo
        \(C\) en función de \(w\) y \(t\). Considera un costo fijo de \(400\)
        mil dólares.
        \item[\textbf{(B)}] ¿Cuánto costaría la construcción si \(w=4\) y
        \(t=7\) kilómetros?
      \end{itemize}
    \end{mwejercicio}
    \column{0.42\textwidth}
    \centering
    \begin{tikzpicture}[scale=0.42, every node/.style={font=\scriptsize}]
      % referencias
      \draw[dashed] (0,7) -- (10.5,7);
      \draw[dashed] (0,4) -- (10.5,4);
      \draw[dashed] (0,0) -- (10.5,0);
      \draw[dashed] (0,0) -- (0,7);
      % recorrido
      \draw[very thick, mwprimario] (0,7) -- (4,4) -- (7,0) -- (10,0);
      \fill (0,7) circle (2.4pt) node[above left] {\(P\)};
      \fill (4,4) circle (2.4pt) node[above right] {\(Q\)};
      \fill (7,0) circle (2.4pt) node[above right] {\(R\)};
      \fill (10,0) circle (2.4pt) node[above right] {\(S\)};
      % cotas verticales
      \draw[<->] (-0.5,4) -- (-0.5,7) node[midway, left] {\(3\) km};
      \draw[<->] (-0.5,0) -- (-0.5,4) node[midway, left] {\(4\) km};
      % cotas horizontales
      \draw[<->] (0,-0.5) -- (4,-0.5) node[midway, below] {\(w\)};
      \draw[<->] (0,-1.2) -- (7,-1.2) node[midway, below] {\(t\)};
      \draw[<->] (0,-1.9) -- (10,-1.9) node[midway, below] {\(10\) km};
    \end{tikzpicture}
  \end{columns}
  \vfill
  \mwtopanchoron
\end{frame}

% ── Ejercicio 11 ────────────────────────────────────────────
\begin{frame}{Ejercicio 11 (A)}
  \begin{mwejercicio}[Ejercicio 11 — Beneficio con producción múltiple \quad \dificil]
    En cada caso, construye la función general que describe el beneficio total
    \(B\) del producto.

    \textbf{(A)} Las empresas de producción múltiple elaboran más de un tipo de
    artículo. Supóngase que una empresa produce dos artículos \(X\) e \(Y\), con
    cantidades \(q_X\) y \(q_Y\). Los artículos están relacionados técnicamente,
    por lo que el costo total \(C\) depende de \(q_X\) y \(q_Y\) (función de
    costos conjuntos). Halla la forma general del beneficio total \(B\) y
    exprésala considerando las diferencias entre monopolio y competencia
    perfecta. Denota las funciones de ingreso y costo como \(I\) y \(C\).
  \end{mwejercicio}
\end{frame}

\begin{frame}{Ejercicio 11 (B)}
  \begin{mwejercicio}[Ejercicio 11 — Discriminación de precios \quad \dificil]
    \textbf{(B)} Las empresas asignan diferentes precios a un mismo producto en
    dos o más mercados monopólicos (discriminación de precios). El monopolista
    trabaja con dos mercados \(X\) e \(Y\), y la producción total es
    \(q=q_X+q_Y\). La función de costo \(C\) depende de la producción total
    \(q\). Halla la forma general del beneficio total \(B\) en empresas que
    practican discriminación de precios. Usa \(I\) para denotar la función de
    ingresos del productor.
  \end{mwejercicio}
\end{frame}

\begin{frame}{Ejercicio 11 (C)}
  \begin{mwejercicio}[Ejercicio 11 — Producción en dos fábricas \quad \dificil]
    \textbf{(C)} Un monopolista produce un único artículo en dos fábricas
    separadas \(X\) e \(Y\). Si \(q_X\) y \(q_Y\) son las cantidades producidas
    en cada fábrica, la producción total es \(q=q_X+q_Y\). El producto es único
    y se vende en el mismo mercado, por lo que la función de ingreso \(I\)
    depende de \(q\). Halla la forma general del beneficio total \(B\) en
    empresas monopólicas que producen en dos fábricas. Usa \(C\) para denotar
    la función de costo del productor.
  \end{mwejercicio}
\end{frame}

% ── Ejercicio 17 ────────────────────────────────────────────
\begin{frame}{Ejercicio 17 (A)}
  \begin{mwejercicio}[Ejercicio 17 — Gorras del Mundial \quad \dificil]
    Con motivo del Mundial de Fútbol, una empresa confeccionará gorras para
    venderlas en Brasilia y Fortaleza. Un estudio determinó:
    \begin{itemize}
      \item En Brasilia se venden \(102\) gorras diarias a \(6\) \$ cada una;
            en Fortaleza, \(51\) gorras diarias a \(8\) \$ la unidad.
      \item Si el precio en Brasilia aumenta \(100\%\) y en Fortaleza disminuye
            \(50\%\), solo se dejaría de vender una gorra en Brasilia y se
            venderían \(2\) más por día en Fortaleza.
      \item Cada gorra cuesta \(3\) \$ y los costos fijos diarios son \(241\) \$.
    \end{itemize}
    \textbf{(A)} Halla una expresión para el beneficio \(B\) diario en términos
    de los precios unitarios \(p_1\) y \(p_2\) en Brasilia y Fortaleza.
    Supón que la demanda en cada ciudad es lineal en el precio.
  \end{mwejercicio}
\end{frame}

\begin{frame}{Ejercicio 17 (B)}
  \begin{mwejercicio}[Ejercicio 17 — Interpretación \quad \medio]
    Con la función de beneficio obtenida en la parte (A):
    \textbf{(B)} calcula \(B(6,8)\) e interpreta el resultado obtenido.
  \end{mwejercicio}
\end{frame}

% ── Ejercicio 20 ────────────────────────────────────────────
\begin{frame}{Ejercicio 20 (A)}
  \begin{mwejercicio}[Ejercicio 20 — Precio y publicidad \quad \dificil]
    Un empresario sabe que si \(P\) es el precio (en dólares) y \(A\) la
    inversión en publicidad (en miles de dólares), la cantidad vendida \(X\) de
    un bien es
    \[
      X=f(P,A)=\frac{32A}{P\,(A^{2}+1)}.
    \]
    Cada unidad tiene un costo unitario de producción de \(2\) \$ y los costos
    totales son la suma de los costos de producción y de la inversión en
    publicidad.

    \textbf{(A)} Determina la función de beneficio \(B\) de la empresa en
    términos de \(A\) y \(P\).
  \end{mwejercicio}
\end{frame}

\begin{frame}{Ejercicio 20 (B) y (C)}
  \begin{mwejercicio}[Ejercicio 20 — Composición y evaluación \quad \dificil]
    \textbf{(B)} Suponiendo que el precio y la inversión en publicidad dependen
    del tiempo \(T\) transcurrido desde la colocación del producto en el
    mercado, de forma que
    \[
      P=\sqrt{T},\qquad A=T-1,
    \]
    determina la función de beneficio de la empresa en términos de \(T\).

    \textbf{(C)} Calcula el beneficio si cada unidad se vende a \(4\) \$ después
    de \(4\) meses desde que el producto comenzó a venderse.
  \end{mwejercicio}
\end{frame}

% ── Ejercicio 21 ────────────────────────────────────────────
\begin{frame}{Ejercicio 21}
  \mwtopanchoroff
  \begin{columns}[T]
    \column{0.50\textwidth}
    \begin{mwejercicio}[Ejercicio 21 — Puertos \quad \muyDificil]
      \small
      Se construirá la infraestructura de transporte entre dos poblados
      aislados: un puerto en el poblado \(A\) (costa de una isla) y otro en el
      litoral continental, más una carretera entre este puerto y el poblado
      \(B\).
      \begin{itemize}
        \item Carretera \(B\)--puerto: \(500\) mil \$ por kilómetro.
        \item Transporte desde \(B\): \(30\) \$ por tonelada-kilómetro por
        carretera y \(12\) \$ por tonelada-kilómetro por vía marítima.
        \item Si la obra excede los \(5\) años: multa de \(2\) millones de \$
        por año de retraso.
      \end{itemize}
    \end{mwejercicio}
    \column{0.45\textwidth}
    \centering
    \begin{tikzpicture}[scale=0.30, every node/.style={font=\scriptsize}]
      % litoral
      \draw[dash pattern=on 5pt off 2pt on 1pt off 2pt]
        (5,-1) -- (5,10.2) node[above] {LITORAL};
      % nivel de B
      \draw[dotted] (0,0) -- (15,0);
      % carretera y tramo maritimo
      \draw[very thick, mwprimario] (0,0) -- (5,4);
      \draw[very thick, mwprimario] (5,4) -- (14,9);
      % puntos
      \fill (0,0) circle (2.4pt) node[below left] {\(B\)};
      \fill (5,0) circle (2.4pt) node[below right] {\(E\)};
      \fill (5,4) circle (2.4pt) node[below right] {\(D\)};
      \fill (14,9) circle (2.4pt) node[left] {\(A\)};
      \node[above left] at (5,4) {PUERTO};
      \node[above right] at (14,9) {PUERTO};
      % cotas
      \draw[<->] (0,-1.1) -- (5,-1.1) node[midway, below] {\(X\) km};
      \draw[<->] (5.6,0) -- (5.6,4) node[midway, right] {\(Y\) km};
      \draw[<->] (5,9.7) -- (14,9.7) node[midway, above] {\(30\) km};
      \draw[<->] (14.8,0) -- (14.8,9) node[midway, right] {\(44\) km};
    \end{tikzpicture}
  \end{columns}
  \vfill
  \mwtopanchoron
\end{frame}

\begin{frame}{Ejercicio 21 (A)}
  \begin{mwejercicio}[Ejercicio 21 — Función de costo total \quad \dificil]
    \textbf{(A)} Determina la función de costo total \(C(x,y,q,t)\), en miles de
    dólares, de construir la infraestructura de transporte descrita, siendo:
    \begin{itemize}
      \item \(x\): distancia horizontal \(BE\) (km);
      \item \(y\): distancia vertical \(ED\) (km);
      \item \(q\): toneladas transportadas por kilómetro recorrido;
      \item \(t\): cantidad de años que dura la construcción.
    \end{itemize}
  \end{mwejercicio}
\end{frame}

\begin{frame}{Ejercicio 21 (B)}
  \begin{mwejercicio}[Ejercicio 21 — Evaluación \quad \medio]
    \textbf{(B)} Calcula \(C(3,4,100,6)\) e interpreta el resultado obtenido.
  \end{mwejercicio}
\end{frame}

% ── Ejercicio 22 ────────────────────────────────────────────
\begin{frame}{Ejercicio 22}
  \begin{mwejercicio}[Ejercicio 22 — Fondo editorial \quad \medio]
    Un fondo editorial vende libros a las librerías de una pequeña ciudad
    europea. Un equipo de \(10\) personas vende enciclopedias y manuales
    técnicos. Cada vendedor tiene un salario anual de \(7000\) euros más una
    comisión de \(30\) euros por enciclopedia y \(35\) euros por manual vendido
    en el año. Si las ventas anuales conjuntas superan las \(900\) unidades, el
    equipo recibe un bono de estímulo \(S\), repartido entre los vendedores en
    proporción a sus ventas. El precio de venta es \(130\) euros por
    enciclopedia y \(140\) euros por manual técnico.

    \textbf{(A)} Determina la función de beneficio \(B\) del fondo editorial en
    función del bono anual \(S\) y de las cantidades \(q_1\) (enciclopedias) y
    \(q_2\) (manuales) vendidas por el equipo.

    \textbf{(B)} ¿Cuál fue el beneficio el año pasado si solo se vendieron
    \(150\) enciclopedias y \(200\) manuales? Interpreta el resultado.
  \end{mwejercicio}
\end{frame}

% ── Ejercicio 26 ────────────────────────────────────────────
\begin{frame}{Ejercicio 26 (A)–(C)}
  \begin{mwejercicio}[Ejercicio 26 — Demanda de cerveza \quad \medio]
    La demanda de cerveza de un consumidor es
    \[
      D(r,p)=\frac{15r}{p}-r^{2}\quad(\text{litros por mes}),
    \]
    donde \(r\) es la renta mensual (en miles de euros) y \(p\) el precio de un
    litro de cerveza. Actualmente \(p=2\) y \(r=2\).
    \begin{itemize}
      \item[\textbf{(A)}] Explica qué significa la notación \(D(r,p)\).
      \item[\textbf{(B)}] Calcula \(D(2,2)\) y explica su significado.
      \item[\textbf{(C)}] ¿Cuál sería la demanda si el consumidor tiene un
      aumento en su renta mensual tal que \(r=2{,}5\)?
    \end{itemize}
  \end{mwejercicio}
\end{frame}

\begin{frame}{Ejercicio 26 (D)–(G)}
  \begin{mwejercicio}[Ejercicio 26 — Incrementos parciales \quad \dificil]
    \begin{itemize}
      \item[\textbf{(D)}] Si además del aumento en la renta del ítem anterior,
      el precio del litro de cerveza sube \(0{,}5\) euros, ¿cuál sería la
      demanda?
      \item[\textbf{(E)}] Expresa los resultados de los dos últimos ítems en
      términos de incrementos parciales de la función de demanda \(D\).
      \item[\textbf{(F)}] Calcula el incremento total de la demanda si
      \(\Delta r=0{,}5\) y \(\Delta p=1\). Interpreta el resultado.
      \item[\textbf{(G)}] ¿Qué signo pareciera tener \(\Delta_r D(3,3)\) si
      \(\Delta r=0{,}5\)? Comprueba si tu conjetura es cierta.
    \end{itemize}
  \end{mwejercicio}
\end{frame}

\begin{frame}{Ejercicio 26 (H)–(J)}
  \begin{mwejercicio}[Ejercicio 26 — Análisis de la demanda \quad \muyDificil]
    \begin{itemize}
      \item[\textbf{(H)}] Con el precio actual, ¿qué renta mensual haría que el
      consumidor limitara su consumo de cerveza a \(9\) litros mensuales?
      \item[\textbf{(I)}] ¿A partir de qué precio el consumidor no estará
      dispuesto a comprar cerveza con su renta actual?
      \item[\textbf{(J)}] ¿Cuánto tendría que reducirse el precio de la cerveza
      para que el consumidor mantuviera su demanda tras haber sufrido un
      recorte salarial de un \(5\%\)?
    \end{itemize}
  \end{mwejercicio}
\end{frame}

% ── Ejercicio 27 ────────────────────────────────────────────
\begin{frame}{Ejercicio 27 (A)}
  \begin{mwejercicio}[Ejercicio 27 — Incrementos parciales \quad \medio]
    Una empresa que fabrica dos productos \(A\) y \(B\) tiene la función de
    beneficio
    \[
      B(x,y)=x^{2}+3y^{2}-xy-20,
    \]
    donde \(x\) e \(y\) son las cantidades producidas de \(A\) y \(B\).
    \textbf{(A)} Halla las funciones de incrementos parciales
    \(\Delta_x B(x,y)\) y \(\Delta_y B(x,y)\).
  \end{mwejercicio}
\end{frame}

\begin{frame}{Ejercicio 27 (B) y (C)}
  \begin{mwejercicio}[Ejercicio 27 — Aplicación \quad \medio]
    \textbf{(B)} La producción actual es \((x,y)=(200,150)\). ¿Cuáles son las
    funciones de incrementos parciales para estos niveles de producción?

    \textbf{(C)} La empresa puede incrementar en \(2\) unidades la producción de
    \(A\) o la de \(B\). ¿Cuál de las dos opciones es más conveniente? ¿Por qué?
  \end{mwejercicio}
\end{frame}

% ── Ejercicio 28 ────────────────────────────────────────────
\begin{frame}{Ejercicio 28 (A) y (B)}
  \begin{mwejercicio}[Ejercicio 28 — Costo en el tiempo \quad \medio]
    La función de costos de una empresa en el instante \(t\) (en años) es
    \[
      C(x,y,t)=(100+200x+10y)\,e^{0{,}01t},
    \]
    donde \(x\) e \(y\) son las cantidades producidas de dos artículos. \(t=0\)
    es el año actual. La producción actual es de \(50\) y \(30\) unidades.
    \begin{itemize}
      \item[\textbf{(A)}] Halla el dominio de la función \(C\).
      \item[\textbf{(B)}] Suponiendo que la producción no se altera de un año al
      siguiente, ¿en qué porcentaje aumentará el costo de producción anualmente?
    \end{itemize}
  \end{mwejercicio}
\end{frame}

\begin{frame}{Ejercicio 28 (C)}
  \begin{mwejercicio}[Ejercicio 28 — Incremento total \quad \dificil]
    \textbf{(C)} Si se prevé que la producción del próximo año será de \(51\) y
    \(32\) unidades de cada artículo, calcula \(\Delta C\) para este período e
    interpreta el resultado obtenido.
  \end{mwejercicio}
\end{frame}

% ── Ejercicio 31 ────────────────────────────────────────────
\begin{frame}{Ejercicio 31 (A)}
  \begin{mwejercicio}[Ejercicio 31 — Propaganda y beneficios \quad \dificil]
    Un empresario gasta diariamente \(S\) dólares en propaganda para promover
    sus productos \(X\) e \(Y\). Las ventas diarias dependen de ese gasto y de
    los precios unitarios \(p_X\) y \(p_Y\):
    \[
      q_X(p_X,p_Y,S)=\frac{3S}{p_X+5}-2p_Y,
      \qquad
      q_Y(p_X,p_Y,S)=\frac{5S}{p_Y+8}-4p_X.
    \]
    Si se producen \(200\) unidades de \(X\) y \(300\) de \(Y\), el costo total
    es de \(2500\) \$ diarios. Los costos unitarios son \(5\) \$ para \(X\) y
    \(3\) \$ para \(Y\). Observa que
    \(C(p_X,p_Y)=c_Xq_X+c_Yq_Y+C_F\).

    \textbf{(A)} Halla una expresión para el beneficio diario \(B\) en términos
    de \(p_X\), \(p_Y\) y \(S\).
  \end{mwejercicio}
\end{frame}

\begin{frame}{Ejercicio 31 (B)–(D)}
  \begin{mwejercicio}[Ejercicio 31 — Estimaciones \quad \dificil]
    \begin{itemize}
      \item[\textbf{(B)}] ¿Cuál sería el beneficio diario si \(p_X=15\) \$,
      \(p_Y=12\) \$ y el gasto diario en propaganda es \(S=1000\) \$?
      \item[\textbf{(C)}] ¿Cuál sería el efecto sobre el beneficio neto diario
      si se mantiene \(p_X=15\) \$, se aumenta \(p_Y\) en \(5\) \$ y
      \(S=1000\) \$?
      \item[\textbf{(D)}] Elabora una lista de todas las ecuaciones que
      relacionan las variables del modelo. Para cada variable indica: nombre,
      notación, tipo (endógena o exógena), descripción y unidades.
    \end{itemize}
  \end{mwejercicio}
\end{frame}

% ── Ejercicio 33 ────────────────────────────────────────────
\begin{frame}{Ejercicio 33}
  \begin{mwejercicio}[Ejercicio 33 — Demanda lineal y beneficio \quad \muyDificil]
    Un comerciante produce dos artículos \(X\) e \(Y\) con precios \(p_X\) y
    \(p_Y\). Cuando vende a \(15\) y \(40\) bolívares por unidad, logra vender
    \(150\) unidades de \(X\) y \(300\) de \(Y\). Cada aumento de \(2\) Bs en el
    precio de cada artículo produce un impacto de \(10\) unidades en \(X\) y de
    \(4\) unidades en \(Y\). El costo total de producir esas cantidades es de
    \(6000\) Bs diarios; los costos fijos son \(2000\) Bs y producir una unidad
    de \(X\) cuesta \(8\) Bs. Determina el beneficio \(B\) diario como función de
    \(p_X\) y \(p_Y\). Usa que la demanda de cada artículo es lineal en su precio.
  \end{mwejercicio}
\end{frame}

% ── Ejercicio 34 ────────────────────────────────────────────
\begin{frame}{Ejercicio 34 (A)}
  \begin{mwejercicio}[Ejercicio 34 — Circo \quad \dificil]
    Llegará a la ciudad un circo. Se estima que asistirán \(x\) personas a cada
    función cuando el precio de la entrada sea \(m\) dólares. Se venderá un
    combo de chucherías cuyo precio es \(p\), con demanda
    \[
      q=h(p)=20\sqrt{29-p}.
    \]
    Cada espectador cuesta \(5\) \$; cada combo cuesta \(1\) \$ de elaboración;
    y hay un costo fijo de \(820\) \$ por función.

    \textbf{(A)} Determina la función de beneficio \(B\) en cada presentación en
    términos de \(x\), \(m\) y \(p\).
  \end{mwejercicio}
\end{frame}

\begin{frame}{Ejercicio 34 (B) y (C)}
  \begin{mwejercicio}[Ejercicio 34 — Evaluación e incrementos \quad \muyDificil]
    \textbf{(B)} Halla el beneficio si asisten \(500\) personas a cada función,
    el precio de la entrada es \(10\) \$ y cada combo se vende a \(4\) \$.

    \textbf{(C)} Por la crisis, se estima que el número de espectadores puede
    ser menor en un \(40\%\). Si se desea obtener el mismo beneficio del ítem
    anterior, ¿cuál debería ser el precio de la entrada si el combo se sigue
    vendiendo a \(4\) \$? Justifica y resuelve usando incrementos. ¿De qué otra
    forma se puede resolver?
  \end{mwejercicio}
\end{frame}

\end{document}
